% The minimal domain model: spawn runs a function on another core,
% join waits for it. Used in M10-L03.
\documentclass[tikz,border=4pt]{standalone}
\usepackage{tikz}
\input{_m10-styles.texinput}
\begin{document}
\begin{tikzpicture}[m10, x=1mm, y=-1mm, font=\large]
  % main domain timeline
  \draw[flow, line width=1.2pt] (0,10) -- (130,10);
  \node[note, anchor=east, font=\large] at (-3,10) {main};

  % spawned domain timeline
  \draw[flow, line width=1.2pt] (44,34) -- (96,34);
  \node[note, anchor=east, font=\large] at (41,34) {Domain B};

  % spawn
  \draw[goodflow, line width=1pt] (44,10) -- (44,34);
  \node[goodnote, anchor=south, font=\large] at (44,5) {\texttt{Domain.spawn f}};
  % work marker
  \node[heapcell, minimum width=26mm, minimum height=9mm, font=\large] at (70,34) {runs \texttt{f ()}};
  % join
  \draw[goodflow, line width=1pt] (96,34) -- (96,10);
  \node[goodnote, anchor=south, font=\large] at (104,5) {\texttt{Domain.join} waits};

  \node[note, anchor=north west, font=\large\itshape] at (0,46)
    {Two domains run at the same time, on different cores.};
\end{tikzpicture}
\end{document}
